\documentclass[12pt,a4paper]{article}

\newif\ifanswers

\ifdefined\PROF
    \answerstrue
\else
    \answersfalse
\fi

\ifdefined\CLASSE
\else
    \def\CLASSE{}
\fi



% Configuration
% ============================================================
% PACKAGES
% ============================================================

% Encodage et langue
\usepackage[utf8]{inputenc}
\usepackage[T1]{fontenc}
\usepackage[french]{babel}

% Mise en page
\usepackage[a4paper, margin=2cm]{geometry}

% Mathématiques
\usepackage{amsmath}
\usepackage{amssymb}

% Dessins
\usepackage{tikz}

% Listes
\usepackage{enumitem}

% Boîtes
\usepackage[most]{tcolorbox}

% Images
\usepackage{graphicx}

% Mise en page avancée
\usepackage{fancyhdr}
\usepackage{multicol}

% Tableaux
\usepackage{booktabs}
\usepackage{tabularx}
\usepackage{colortbl}

\usepackage[normalem]{ulem}

\usepackage{pgfplots}
\pgfplotsset{compat=1.18}
\usepackage{pgf-pie}


\usepackage{ifthen}

% ============================================================
% COULEURS
% ============================================================

% ------------------------------------------------------------
% Couleurs principales
% ------------------------------------------------------------

\definecolor{Bleu}{RGB}{40, 100, 180}
\definecolor{Vert}{RGB}{50, 150, 90}
\definecolor{Orange}{RGB}{230, 140, 40}

\definecolor{BleuClair}{RGB}{235, 245, 255}
\definecolor{VertClair}{RGB}{235, 250, 240}
\definecolor{OrangeClair}{RGB}{255, 245, 225}

% ------------------------------------------------------------
% Couleurs complémentaires
% ------------------------------------------------------------

% Rappels de cours — violet
\definecolor{Violet}{RGB}{125, 85, 170}
\definecolor{VioletClair}{RGB}{245, 238, 252}

% Théorèmes — rouge
\definecolor{Rouge}{RGB}{190, 55, 55}
\definecolor{RougeClair}{RGB}{255, 238, 238}

% Lemme — indigo
\definecolor{Indigo}{RGB}{75, 85, 155}
\definecolor{IndigoClair}{RGB}{240, 241, 252}

% Points méthode — turquoise
\definecolor{Turquoise}{RGB}{35, 145, 150}
\definecolor{TurquoiseClair}{RGB}{235, 250, 250}

% Erreurs / attention — jaune
\definecolor{Jaune}{RGB}{210, 165, 35}
\definecolor{JauneClair}{RGB}{255, 250, 225}

% ------------------------------------------------------------
% Neutres
% ------------------------------------------------------------

\definecolor{GrisClair}{RGB}{245, 245, 245}
\definecolor{Gris}{RGB}{100, 100, 100}


% ============================================================
% COMMANDES PERSONNALISÉES
% ============================================================

\newcommand{\bb}[1]{\textbf{#1}}
\newcommand{\ii}[1]{\textit{#1}}

\newcommand{\cm}{\,\mathrm{cm}}
\newcommand{\m}{\,\mathrm{m}}

\newcommand{\reponse}[1]{\boxed{#1}}

% ============================================================
% TABLEAUX
% ============================================================

\newcommand{\ligneentete}{%
    \rowcolor{Bleu}%
    \color{white}\bfseries
}

\newcommand{\entete}[1]{%
    \cellcolor{Bleu}\color{white}\bb{#1}%
}


\NewDocumentCommand{\hl}{O{JauneClair}m}{%
  \tcbox[
    on line,
    colback=#1,
    colframe=#1,
    boxrule=0pt,
    arc=4pt,
    left=0pt,
    right=0pt,
    top=1pt,
    bottom=1pt
  ]{#2}%
}

\NewDocumentCommand{\hlb}{O{JauneClair}m}{%
  \tcbox[
    on line,
    colback=#1,
    colframe=#1,
    boxrule=0pt,
    arc=4pt,
    left=0pt,
    right=0pt,
    top=1pt,
    bottom=1pt
    ]{\bb{#2}}%
}
\newlength{\answerwidth}


\newcommand{\answer}[2][VertClair]{%
    \ifmmode
        \settowidth{\answerwidth}{$\displaystyle #2$}%
        \ifanswers
            {\setlength{\fboxsep}{0pt}%
             \colorbox{#1}{$\displaystyle #2$}}%
        \else
            \makebox[0pt][l]{%
                \makebox[\answerwidth][l]{%
                    \leaders\hbox to 0.5em{\hss.\hss}\hfill
                }%
            }%
            \phantom{#2}%
        \fi
    \else
        \settowidth{\answerwidth}{#2}%
        \ifanswers
            \hl[#1]{#2}%
        \else
            \makebox[\answerwidth][l]{%
                \leaders\hbox to 0.5em{\hss.\hss}\hfill
            }%
        \fi
    \fi
}


\newcommand{\graphanswer}[1]{%
    \ifanswers
        #1%
    \else
        0%
    \fi
}

\newcommand{\exo}[1]{%
    \includegraphics[
        width=\linewidth,
        keepaspectratio
    ]{#1}%
}



\newcounter{exercice}
\newtcolorbox{exercicebox}[1][\unskip]{
    colback=BleuClair,
    colframe=Bleu,
    boxrule=1pt,
    arc=3mm,
    title=\bb{Exercice \stepcounter{exercice}\theexercice #1},
    left=5mm,
    right=5mm,
    top=2mm,
    bottom=2mm
}


\newcounter{probleme}
\newtcolorbox{probbox}[1][\unskip]{
    colback=RougeClair,
    colframe=Rouge,
    boxrule=1pt,
    arc=3mm,
    title=\bb{Problème \stepcounter{probleme}\theprobleme #1},
    left=5mm,
    right=5mm,
    top=2mm,
    bottom=2mm
}

\newenvironment{tableau}[1]
{%
    \renewcommand{\arraystretch}{1.3}%
    \begin{center}%
    \begin{tabularx}{\linewidth}{#1}%
    \toprule
}
{%
    \bottomrule
    \end{tabularx}%
    \end{center}%
    \renewcommand{\arraystretch}{1}%
}




\pagestyle{fancy}
\fancyhf{}

% ===============================================
% ###### INFO ENTETE ########
% ===============================================
\def\titre{Priorités opératoires et nombres relatifs}
\def\soustitre{}

% En-tête
\fancyhead[L]{Nom: \hspace{4cm} Prénom: \hspace{4cm}Classe:}

% Numéro de page
\fancyfoot[C]{\thepage}
% Ligne sous l'en-tête
\renewcommand{\headrulewidth}{0.4pt}
% Pas de ligne au-dessus du pied de page
\renewcommand{\footrulewidth}{0pt}


\begin{document}
\renewcommand{\labelitemi}{\textbullet}

\begin{center}
    {\Large\bb{Devoir surveillé 1}}
\end{center}



\begin{exercicebox}[\it{(Priorités opératoires)} \hspace*{\fill}\bb{4pts}]
    Effectuer les calculs suivants \bb{en détaillant les étapes}
    \begin{minipage}{0.49\textwidth}
        \begin{align*}
        &\text{A = }2+5\times4\\
        &\text{B = }10-4+10\times2
        \end{align*}
    \end{minipage}
    \hfill
    \begin{minipage}{0.49\textwidth}
        \begin{align*}
            &\text{C = }7\times[2\times(8-3)]\\
            &\text{D = } 2\times3\div4\times(18-16)
        \end{align*}
    \end{minipage}
\end{exercicebox}

\begin{exercicebox}[ \it{(Opérations de base sur les nombres
    relatifs)}\hspace*{\fill}\bb{3pts}]
    Effectuer les calculs suivants \bb{en détaillant les étapes}

    \begin{minipage}{0.49\textwidth}
    \begin{align*}
        &\text{A = }(-10)+(-2)\\
        &\text{B = }(-6)+(+14)\\
        &\text{C = }(-18)-(-8)
    \end{align*}
    \end{minipage}
    \hfill
    \begin{minipage}{0.49\textwidth}
    \begin{align*}
        &\text{D = }13-(+4)\\
        &\text{E = }2\times(-4)\\
        &\text{F = }(+24)\div(-4)
    \end{align*}
    \end{minipage}
\end{exercicebox}

\begin{exercicebox}[ \it{(Enchaîner les calculs avec des nombres relatifs)}
    \hspace*{\fill}\bb{3.5pts}]
    \begin{enumerate}
        \item Déterminer, sans calculer le résultat, si les résultats des 
            calculs suivants sont positifs ou négatifs. \bb{Justifier par une phrase.}: \\
            $$\text{a. } 4\times(-12) \hspace{1cm}\text{b. }
            (-34)\times(-2)\times(+56) \hspace{1cm} \text{c. }
            18\div(-3)\div(+4)$$
        \item Effectuer les calculs suivants \bb{en détaillant les étapes}:
    \begin{align*}
        &\text{A = }18-(-6)+(-4)\\
        &\text{B = }18\div(+6)\times(-4)
    \end{align*}
    \end{enumerate}
\end{exercicebox}

\begin{exercicebox}[ \it{(Priorités opératoires avec des nombres
    relatifs)} \hfill \bb{6pts}]
    Effectuer les calculs suivants \bb{en détaillant les étapes}
    \begin{align*}
        &\text{A = }9\times[\hspace{2pt}(-3)\times(-2)+5]\\
        &\text{B = }[18+(-9)]\times[(55-50)+(-6)]\\
        &\text{C = }[-2\times(4+(-6))]\div[5\times11+(-53)]
    \end{align*}
\end{exercicebox}

\begin{probbox}[\hspace*{\fill}3.5pts]
    Zahra se trouvait au rez-de-chaussée d'une tour de 31 étages.
    \newline
    \newline
        Elle est entrée dans l'ascenseur et
        a appuyé sur le bouton 3, mais l'ascenseur est
        devenu fou : il est descendu de trois étages,
        monté de cinq étages, descendu de quatre
        étages, monté d'un étage. Et tout cela trois fois
        de suite, avant de libérer Zahra !
    \newline
    \newline
    Lorsqu'elle est sortie de l'ascenseur, était-elle à l'étage désiré ?
    \newline
    \newline
    Détailler le raisonnement.

\end{probbox}
\end{document}
